Each step gives the setting we used and what it changed (Table~\ref{tab:key}), with every method at one default setting; Table~\ref{tab:cost} gives the training cost.
\begin{enumerate}\setlength{\itemsep}{1pt}
  \item \textbf{Classification: adapt the encoder's attention weights.} Rank-16 LoRA ($\alpha=32$, learning rate $10^{-4}$) on every block added $\numCiClsloralinAllMeandiff$ AUROC over the linear probe and $\numCiClsloraattnAllMeandiff$ over the attentive probe, positive in all \numCiClsloraattnAllClustersN{} cohort--backbone pairs. Fine-tuning the attention weights in full was ahead of LoRA in \numSeedsclsAttnftLoraAllPosN{} of \numSeedsclsAttnftLoraAllPairsN{} seed pairs, by $\numCiClsattnftloraAllMeandiff$ on average, at similar GPU memory and time per step; LoRA stores \numParamsratioPillarZAttnftLora{} to \numParamsratioSpectreAttnftLora{} times fewer weights per model. Fine-tuning the MLP weights instead was close on TAP-CT and Pillar-0 and behind on SPECTRE. With LoRA, use plain LoRA: no variant improved on it by more than $\numVsloraAllBabmilMeandiff$.
  \item \textbf{Classification: start with one model per cohort.} It was on average level with per-label models ($\numCiClsjointAllMeandiff$) at $1/K$ of the training, but per cohort the difference ranged from $\numJointdeltaMin$ to $\numJointdeltaMax$, so compare the two on selection data when the budget allows.
  \item \textbf{Segmentation: put a multi-depth decoder on the frozen encoder.} It added $\numSegdeltaTapctLinDecMeandiff$ (TAP-CT) and $\numSegdeltaSpectreLinDecMeandiff$ (SPECTRE) Dice over a linear probe, on every dataset. Adapting the encoder beneath it changed Dice by at most $\numSegdeltaSpectreDecDoraMeandiff$ on average (DoRA on SPECTRE); full fine-tuning ranked best but lost more than 0.1 Dice on some small structures, so check Dice per structure.
  \item \textbf{Segmentation: train one model per structure when the budget allows.} On both 15-structure datasets, per-structure models at 8000 steps were ahead of the joint model for the frozen decoder, LoRA and full fine-tuning, and at matched total compute per-structure decoders were still ahead of the joint decoder in three of four pairs. Otherwise check the Dice of small structures, where the joint models failed.
  \item \textbf{Full fine-tuning: lower the learning rate of early blocks.} With the encoder learning rate of the other methods it did not train on SPECTRE (near-constant predictions); layer-wise decay of 0.75 per block fixed it. Check that held-out predictions vary between cases.
\end{enumerate}

\paragraph{Limitations.} Each method has one untuned default setting, and only full fine-tuning was retrained after it failed, so the ranking among the stronger methods may change with tuning; seeds were repeated only for joint models. Per-label tests have little power: LoRA was significantly ahead of the linear probe on \numSigbysizeSmallSig\% of labels with fewer than 100 held-out cases and on \numSigbysizeLargeSig\% of those with 400 or more. Segmentation had not converged and is scored on each backbone's working grid, so it ranks adaptation methods rather than measuring the distance to nnU-Net \citep{isensee2021nnunet} or to models pretrained on segmentation labels \citep{li2024suprem,he2024vista3d}. SPECTRE saw AMOS, LUNA25, LiTS, MSD Pancreas and most held-out KiTS23 cases, without labels, in self-supervised pretraining (Table~\ref{tab:overlap}), which may favor it over the other backbones on these sets.
